%! Solving Quadratics by Factoring — Unit 2, Lesson 2.2 — Slides (Beamer)
% SLIDES-FIRST shape (course shape 2026-09-25; this lesson is its pilot). The deck IS the lesson;
% its printed handout (lesson02_slides.pdf, 3-up with a notes column) is the student's notes.
%   title → targets → warm-up → 4 × { definition → example → now you try } → wrap-up → homework
% One anchor runs through it: x^2-2x-3 = (x+1)(x-3), the quadratic of 2.0 and 2.1, now factored by
% hand. Cycle 4's Now you try is the crux (dividing by x loses x = 0; the product must be ZERO).
% Answers only inside \reveal{} — the handout is a Beamer handout-mode compile and drops them.
\documentclass[aspectratio=169, 11pt]{beamer}
\usepackage{algebra2-beamer}   % \forestheader{}, \sectionlabel[color]{}, \reveal[n]{}
% Coordinate grid + axes: {xmin}{xmax}{ymin}{ymax}
\newcommand{\lgrid}[4]{%
  \draw[step=1, linegray!45, thin] (#1,#3) grid (#2,#4);
  \draw[->, thick, charcoal] (#1,0)--(#2,0) node[below right, font=\tiny]{$x$};
  \draw[->, thick, charcoal] (0,#3)--(0,#4) node[left, font=\tiny]{$y$};}
\setbeamercolor{block title}{bg=forest, fg=white}
\setbeamercolor{block body}{bg=forestbg, fg=charcoal}

\begin{document}

% TITLE SLIDE (CourseName is not defined in beamer, so the name is literal)
{
\setbeamercolor{background canvas}{bg=forest}
\begin{frame}[plain]
  \vspace{2.8cm}\hspace{0.55cm}%
  \begin{minipage}{0.9\textwidth}
    {\color{goldacc}\bfseries\small LESSON 2.2}\\[0.5cm]
    {\color{white}\bfseries\LARGE Solving Quadratics by Factoring}\\[1.2cm]
    {\color{forestmid}\small Algebra 2~$\cdot$~Unit 2: Quadratic Functions}
  \end{minipage}
\end{frame}
}

% LEARNING TARGETS
\begin{frame}[t, plain]
  \forestheader{Today's targets}
  \sectionlabel{What you will be able to do}
  \begin{itemize}\setlength{\itemsep}{5pt}
    \item \textbf{Factor} $x^2+bx+c$ by the \textbf{product--sum} search, taking out the
          \textbf{GCF} first.
    \item Factor $ax^2+bx+c$ by \textbf{grouping}.
    \item Recognize a \textbf{difference of squares} and a \textbf{perfect-square trinomial}.
    \item Solve a quadratic equation with the \textbf{Zero Product Property}, and read its
          \textbf{roots} as the graph's $x$-intercepts.
  \end{itemize}
  \vspace{3pt}
  \begin{block}{How today runs}
    Warm-up (5) $\to$ four rounds of definition, example, and \textbf{now you try} (40) $\to$
    wrap-up (5) $\to$ \textbf{start the homework} (10).
  \end{block}
\end{frame}

% WARM-UP — multiply (the move factoring undoes), a product-sum search, and the zeros read from
% 2.1's factored form. The hold-on block leaves the question cycle 1 answers.
\begin{frame}[t, plain]
  \forestheader{Warm-Up}
  \sectionlabel{5 minutes --- in the notes column, alone}
  \begin{enumerate}\setlength{\itemsep}{6pt}
    \item Multiply: $(x+1)(x-3)$. \quad \reveal{\textbf{$x^2-2x-3$}}
    \item Find two integers whose product is $-12$ and whose sum is $1$.
          \quad \reveal{\textbf{$4$ and $-3$}}
    \item Where does $y=(x+1)(x-3)$ cross the $x$-axis? \quad
          \reveal{\textbf{$x=-1$ and $x=3$} --- each factor is $0$ there.}
  \end{enumerate}
  \vspace{4pt}
  \begin{block}{Hold on to this}
    Item 1 turned a product into $x^2-2x-3$. If all you are given is $x^2-2x-3$, how do you
    get the product back?
  \end{block}
\end{frame}

% ── CYCLE 1 — FACTORING x^2+bx+c, GCF FIRST ──────────────────────────────────
\begin{frame}[t, plain]
  \forestheader{Definition: factoring $x^2+bx+c$}
  \sectionlabel{Idea 1 --- un-multiplying}
  \begin{block}{Factoring}
    To \textbf{factor} a polynomial is to write it as a \emph{product}. Always take out the
    \textbf{greatest common factor (GCF)} first.
  \end{block}
  \vspace{4pt}
  \begin{block}{The product--sum search}
    Because $(x+p)(x+q)=x^2+(p+q)x+pq$, factor $x^2+bx+c$ by finding two integers $p$ and $q$
    with \[ p\cdot q=c \qquad\text{and}\qquad p+q=b. \] Then $x^2+bx+c=(x+p)(x+q)$.
  \end{block}
\end{frame}

\begin{frame}[t, plain]
  \forestheader{Example 1: factor $x^2-2x-3$, then $2x^2-4x-6$}
  \sectionlabel{Watch --- I do this one}
  \begin{columns}[T]
    \begin{column}{0.48\textwidth}
      \textbf{(a)} $x^2-2x-3$: product $-3$, sum $-2$.
      \pause
      \begin{center}
      \renewcommand{\arraystretch}{1.15}
      \begin{tabular}{c|c}
        pair & sum \\ \hline
        $1,\,-3$ & $\mathbf{-2}$ \checkmark \\
        $-1,\,3$ & $2$
      \end{tabular}
      \end{center}
      \pause
      $x^2-2x-3=(x+1)(x-3)$. \\[3pt]
      \emph{Check by multiplying:} the warm-up's item 1.
    \end{column}
    \pause
    \begin{column}{0.48\textwidth}
      \textbf{(b)} $2x^2-4x-6$: every term has a $2$.
      \[ 2x^2-4x-6=2\,(x^2-2x-3) \]
      \pause
      \[ =2\,(x+1)(x-3). \]
      \emph{GCF first} --- then the trinomial left inside is the one we just factored.
    \end{column}
  \end{columns}
\end{frame}

\begin{frame}[t, plain]
  \forestheader{Now you try}
  \sectionlabel[goldacc]{In the notes column --- then check}
  Factor completely.
  \begin{enumerate}\setlength{\itemsep}{4pt}
    \item $x^2+5x-14$
    \item $3x^2+3x-18$
  \end{enumerate}
  \vspace{6pt}
  \reveal{%
    \begin{block}{Check}
      1.\ Product $-14$, sum $5$: $7$ and $-2$. \quad $(x+7)(x-2)$ \\[2pt]
      2.\ GCF $3$ first: $3(x^2+x-6)$; product $-6$, sum $1$: $3$ and $-2$. \quad
      $3(x+3)(x-2)$
    \end{block}}
\end{frame}

% ── CYCLE 2 — ax^2+bx+c BY GROUPING ──────────────────────────────────────────
\begin{frame}[t, plain]
  \forestheader{Definition: factoring $ax^2+bx+c$ by grouping}
  \sectionlabel{Idea 2 --- when $a\ne1$}
  \begin{block}{Grouping (the $a\cdot c$ method)}
    To factor $ax^2+bx+c$ with $a\ne1$ (after any GCF):
    \begin{enumerate}\setlength{\itemsep}{2pt}
      \item Find two integers whose product is $\mathbf{a\cdot c}$ and whose sum is $\mathbf{b}$.
      \item Use them to \textbf{split} the middle term $bx$ into two terms.
      \item \textbf{Group} the four terms in pairs, take the GCF out of each pair, and take out
            the common binomial.
    \end{enumerate}
  \end{block}
\end{frame}

\begin{frame}[t, plain]
  \forestheader{Example 2: factor $2x^2+5x-3$}
  \sectionlabel{Watch --- I do this one}
  $a\cdot c=2(-3)=-6$, \; $b=5$: \; the pair is $6$ and $-1$ \;($6\cdot(-1)=-6$, $6+(-1)=5$).
  \pause
  \[
  \begin{aligned}
    2x^2+5x-3 &= 2x^2+6x-x-3 && \text{split } 5x \text{ as } 6x-x \\
              &= 2x(x+3)-1(x+3) && \text{GCF of each pair} \\
              &= (2x-1)(x+3)    && \text{take out } (x+3)
  \end{aligned}
  \]
  \pause
  \emph{Check:} $(2x-1)(x+3)=2x^2+6x-x-3=2x^2+5x-3$ \checkmark
\end{frame}

\begin{frame}[t, plain]
  \forestheader{Now you try}
  \sectionlabel[goldacc]{In the notes column --- then check}
  Factor by grouping: \quad $3x^2-10x+8$
  \vspace{10pt}
  \reveal{%
    \begin{block}{Check}
      $a\cdot c=24$, sum $-10$: $-6$ and $-4$. \\[2pt]
      $3x^2-6x-4x+8=3x(x-2)-4(x-2)=(3x-4)(x-2)$ \\[2pt]
      \emph{Watch the sign:} the second pair is $-4x+8$, so its GCF is $-4$, not $4$ ---
      that is what makes both brackets $(x-2)$.
    \end{block}}
\end{frame}

% ── CYCLE 3 — SPECIAL PATTERNS ──────────────────────────────────────────────
\begin{frame}[t, plain]
  \forestheader{Definition: two special patterns}
  \sectionlabel{Idea 3 --- factor on sight}
  \begin{columns}[T]
    \begin{column}{0.58\textwidth}
      \begin{block}{Difference of squares}
        \[ a^2-b^2=(a+b)(a-b) \]
      \end{block}
      \vspace{2pt}
      \begin{block}{Perfect-square trinomial}
        \[ a^2+2ab+b^2=(a+b)^2 \qquad a^2-2ab+b^2=(a-b)^2 \]
      \end{block}
      \vspace{2pt}
      {\small A \emph{sum} of squares, $a^2+b^2$, does \textbf{not} factor over the integers.}
    \end{column}
    \begin{column}{0.38\textwidth}
      \centering
      \begin{tikzpicture}[scale=0.36]
        \lgrid{-1}{7}{-2}{6}
        \begin{scope}\clip (-1,-2) rectangle (7,6);
          \draw[very thick, forest, domain=0.4:5.6, samples=60, smooth]
            plot (\x,{(\x-3)*(\x-3)});
        \end{scope}
        \fill[redacc] (3,0) circle (4pt) node[below, font=\tiny, charcoal]{$3$};
      \end{tikzpicture}\\[-2pt]
      {\scriptsize $y=x^2-6x+9=(x-3)^2$ \emph{touches} the axis once: \\ a \textbf{double root}.}
    \end{column}
  \end{columns}
\end{frame}

\begin{frame}[t, plain]
  \forestheader{Example 3: three on sight}
  \sectionlabel{Watch --- I do this one}
  \begin{itemize}\setlength{\itemsep}{8pt}
    \item $x^2-25$: \; $x^2-5^2$, a difference of squares. \pause \; $=(x+5)(x-5)$
    \pause
    \item $4x^2-9$: \; $(2x)^2-3^2$. \pause \; $=(2x+3)(2x-3)$
    \pause
    \item $x^2-6x+9$: \; first and last are squares ($x^2$, $3^2$), and the middle is
          $-2\cdot x\cdot3$. \pause \; $=(x-3)^2$
  \end{itemize}
  \pause
  \vspace{4pt}
  \emph{Always test:} are the first and last terms perfect squares? Then look at the middle.
\end{frame}

\begin{frame}[t, plain]
  \forestheader{Now you try}
  \sectionlabel[goldacc]{In the notes column --- then check}
  Factor completely, or say it does not factor.
  \begin{enumerate}\setlength{\itemsep}{4pt}
    \item $9x^2-16$ \qquad 2.\ $x^2+10x+25$ \qquad 3.\ $x^2+16$
  \end{enumerate}
  \vspace{6pt}
  \reveal{%
    \begin{block}{Check}
      1.\ $(3x)^2-4^2=(3x+4)(3x-4)$ \\[2pt]
      2.\ $x^2+2\cdot5\cdot x+5^2=(x+5)^2$ \\[2pt]
      3.\ A \emph{sum} of squares --- \textbf{does not factor}. (No two integers have product
      $16$ and sum $0$.)
    \end{block}}
\end{frame}

% ── CYCLE 4 — ZERO PRODUCT PROPERTY (the crux) ──────────────────────────────
\begin{frame}[t, plain]
  \forestheader{Definition: the Zero Product Property}
  \sectionlabel{Idea 4 --- from factoring to solving}
  \begin{columns}[T]
    \begin{column}{0.58\textwidth}
      \begin{block}{Zero Product Property}
        If $A\cdot B=0$, then $A=0$ or $B=0$.
      \end{block}
      \vspace{2pt}
      \begin{block}{Solving by factoring}
        \begin{enumerate}\setlength{\itemsep}{1pt}
          \item Move every term to one side: $\;\ldots=0$.
          \item Factor completely.
          \item Set each factor equal to $0$ and solve.
        \end{enumerate}
        The solutions are the \textbf{roots}; on the graph of $y=ax^2+bx+c$ they are the
        $x$-intercepts (the \textbf{zeros}).
      \end{block}
    \end{column}
    \begin{column}{0.38\textwidth}
      \centering
      \begin{tikzpicture}[scale=0.36]
        \lgrid{-3}{5}{-5}{3}
        \begin{scope}\clip (-3,-5) rectangle (5,3);
          \draw[very thick, forest, domain=-1.5:3.5, samples=70, smooth]
            plot (\x,{(\x+1)*(\x-3)});
        \end{scope}
        \fill[redacc] (-1,0) circle (4pt) node[above left, font=\tiny, charcoal]{$-1$};
        \fill[redacc] (3,0) circle (4pt) node[above right, font=\tiny, charcoal]{$3$};
      \end{tikzpicture}\\[-2pt]
      {\scriptsize $y=x^2-2x-3$: zeros $-1$ and $3$.}
    \end{column}
  \end{columns}
\end{frame}

\begin{frame}[t, plain]
  \forestheader{Example 4: solve $x^2=2x+3$}
  \sectionlabel{Watch --- I do this one}
  \[
  \begin{aligned}
    x^2 &= 2x+3 \\
    x^2-2x-3 &= 0            && \text{everything on one side} \\
    (x+1)(x-3) &= 0          && \text{factor (cycle 1)} \\
    x+1=0 \quad\text{or}\quad x-3 &= 0 && \text{Zero Product Property} \\
    x=-1 \quad\text{or}\quad x &= 3
  \end{aligned}
  \]
  \pause
  \emph{Check:} $(-1)^2=1=2(-1)+3$ \checkmark \quad $3^2=9=2(3)+3$ \checkmark \quad --- and these
  are the graph's $x$-intercepts.
\end{frame}

\begin{frame}[t, plain]
  \forestheader{Now you try}
  \sectionlabel[redacc]{The one that matters --- in the notes column, then check}
  Solve.
  \begin{enumerate}\setlength{\itemsep}{4pt}
    \item $x^2=5x$ \qquad\qquad 2.\ $(x-1)(x+2)=4$
  \end{enumerate}
  \vspace{4pt}
  \reveal{%
    \begin{block}{Check}
      1.\ $x^2-5x=0 \Rightarrow x(x-5)=0 \Rightarrow x=0$ or $x=5$. \;
      \emph{Dividing both sides by $x$ gives only $x=5$ --- it throws away $x=0$.} \\[3pt]
      2.\ The product is $4$, not $0$ --- so ``$x-1=4$ or $x+2=4$'' is wrong. Multiply out:
      $x^2+x-2=4 \Rightarrow x^2+x-6=0 \Rightarrow (x+3)(x-2)=0 \Rightarrow x=-3$ or $x=2$.
    \end{block}}
\end{frame}

% WRAP-UP — the same content as the cover's Keep in Mind.
\begin{frame}[t, plain]
  \forestheader{Wrap-up: what to keep}
  \sectionlabel{5 minutes}
  \begin{itemize}\setlength{\itemsep}{3pt}
    \item \textbf{Factor} = write as a product. \textbf{GCF first}, always.
    \item $x^2+bx+c$: two integers with product $c$, sum $b$. \quad $ax^2+bx+c$: product
          $a\cdot c$, sum $b$, then \textbf{group}.
    \item $a^2-b^2=(a+b)(a-b)$; \; $a^2\pm2ab+b^2=(a\pm b)^2$; \; $a^2+b^2$ does not factor.
    \item \textbf{Zero Product Property:} $AB=0 \Rightarrow A=0$ or $B=0$. The roots are the
          $x$-intercepts.
  \end{itemize}
  \vspace{3pt}
  \begin{block}{Watch out}
    The property is about \textbf{zero}. Get everything on one side $=0$ \emph{before} you factor
    --- and never divide both sides by $x$: that loses the root $x=0$.
  \end{block}
\end{frame}

% HOMEWORK — source: DeltaMath, printed (user choice for 2.2, 2026-09-25).
\begin{frame}[t, plain]
  \forestheader{Homework --- start it now}
  \sectionlabel[goldacc]{10 minutes in class, alone}
  \begin{block}{Homework --- scored, due the first class after two study halls}
    The \textbf{DeltaMath} pages in your packet: factoring and solving by factoring. Start the
    first problem now; the rest is due the first class after two study halls. I will announce
    the date in class and on TurtleNet.
  \end{block}
  \vspace{6pt}
  \textbf{Next:} Lesson 2.3 --- when a quadratic will not factor, solve it with square roots
  and by completing the square.
\end{frame}

\end{document}
